\chapter{Produto Educacional e Sequência Didática}

\section{Concepção do produto educacional}

O produto educacional tem duas partes: uma plataforma interativa em português brasileiro e uma sequência didática para o Ensino Médio. A plataforma funciona no navegador, trabalha com cenários sintéticos e faz os cálculos no próprio computador. Não exige cadastro nem conta de usuário, e não envia respostas nem coleta informações pessoais. Isso facilita o acesso e permite o uso com conectividade limitada, desde que o recurso esteja no computador ou na projeção da aula.

A plataforma apoia a mediação do professor no ensino de Computação. Não substitui a explicação, a conversa nem a intervenção dele.

A proposta tem cinco módulos introdutórios de \textit{machine learning}: regressão linear, regressão logística, \textit{k}-vizinhos mais próximos, árvore de decisão e \textit{k}-means. Eles não formam um catálogo de comandos nem conteúdos soltos. Cada módulo mostra uma tarefa diferente com dados (prever um valor, classificar um caso, formar grupos), e a turma compara o que muda de um para outro. Todos têm lugar central na sequência, ainda que cada um mobilize uma camada matemática própria. A regressão linear traz tendência, função, previsão, resíduo e erro. A logística traz probabilidade e limiar de decisão. No \textit{k}-NN pesa a distância entre pontos; no \textit{k}-means, a média e a posição dos centroides; na árvore, as condições encadeadas.

A diferença de formalização não diminui nenhum módulo. Vem do repertório de estudantes iniciantes e da necessidade de tornar visível, no gráfico ou na regra, uma relação que a turma consiga discutir durante a aula.

Perguntas que voltam em todos os módulos ligam os modelos, mais do que a presença na mesma plataforma. A turma pergunta qual é o dado disponível, de onde ele vem, que saída se espera e qual alteração pode modificá-la. A resposta muda conforme o modelo. Uma reta produz uma previsão numérica; a regressão logística leva a uma probabilidade, que se compara com um limiar; o \textit{k}-NN depende dos exemplos próximos; a árvore percorre condições; o \textit{k}-means reorganiza pontos sem categorias dadas de antemão.

Com as mesmas perguntas em cada módulo, os estudantes comparam os modelos sem reduzi-los a uma lista de nomes, e o professor percebe se a dificuldade está em um conceito específico ou na leitura de dados e gráficos, que atravessa todos os módulos.

\section{Organização pedagógica da plataforma}

Cada módulo começa por uma situação-problema, pela origem dos dados e por uma pergunta sobre o que o modelo tenta aprender. Antes de mexer nos controles, o estudante prevê uma consequência. Depois altera um dado, um parâmetro ou uma regra definidos para o experimento, observa o que muda no gráfico, na classificação ou no agrupamento e explica o que aconteceu.

O ciclo prever, experimentar, observar e explicar se repete em todos os módulos e expõe o raciocínio que acompanha cada manipulação.

Mover a inclinação da reta leva o estudante a comparar a previsão com o valor observado. A troca de um limiar leva a perguntar por que a classe muda mesmo quando a porcentagem calculada continua a mesma. Com o número de vizinhos, as divisões da árvore e a posição dos centroides, a lógica é a mesma.

A sequência abre com recursos de inteligência artificial que os estudantes já usam e chega logo a um problema concreto com dados. A Matemática só entra quando a situação, o gráfico ou a decisão do modelo a exige.

Toda representação visual vem com descrição em texto, tabela ou regra, para que a compreensão não dependa só da imagem. As visualizações em 3D do protótipo complementam e não são necessárias para entender o conceito central. A interface deve preservar o controle por teclado e respeitar a preferência por redução de movimento.

\section{Cenários e dados dos experimentos}

Os dados das atividades serão quase todos sintéticos. Nos encontros 2 e 8, a turma organizará pequenas tabelas com dados simples, voluntários e sem nome, ou com medidas de objetos da sala. Essas tabelas servem ao ensino e ficam fora do corpus da pesquisa. Os demais cenários evitam notas, comportamentos, trajetórias e qualquer informação pessoal dos estudantes.

Os cenários são uma primeira proposta e podem mudar depois que a turma e a escola forem definidas. Na regressão linear, a ideia é relacionar temperatura e consumo fictício de energia em uma sala. Na regressão logística, a leitura de um sensor fictício permite discutir probabilidade e alerta. O \textit{k}-NN poderá classificar objetos ou situações simuladas por duas características mensuráveis, e a árvore de decisão trabalhará uma irrigação simulada, com umidade e temperatura. No \textit{k}-means, pontos de coleta fictícios em um mapa de bairro servem para explorar proximidade, centroides e agrupamentos sem rótulos prévios. Nenhum cenário entra só pelo apelo visual. Cada um precisa fazer a mudança de um dado ou parâmetro produzir uma consequência observável, que a turma consiga discutir com a Matemática da série.

Antes da aplicação, os exemplos serão revistos quanto à clareza da linguagem, ao tempo que a atividade toma e à pertinência para a turma.

\section{Sequência didática provisória}

A sequência foi planejada, por ora, para oito encontros de cerca de cinquenta minutos. O laboratório de informática da escola, com 20 computadores e acesso à internet, já está confirmado. O número de aulas, a turma e o componente curricular dependem ainda da viabilidade do campo. O cenário preferencial é uma turma de Ensino Médio em que o pesquisador é o professor, possivelmente na disciplina Projeto Permanente por Afinidade; a proposta pode migrar para outro componente, inclusive Matemática, se a lotação e o calendário escolar exigirem.

Todo encontro mantém a mesma estrutura, mesmo quando o modelo muda. A aula abre com uma situação ou pergunta que delimita o problema. Cada estudante faz uma previsão individual breve, explora um controle ou uma representação com o colega, a turma discute o que observou e o professor fecha com uma síntese. Isso dá continuidade às aulas e deixa claro que a plataforma não é um intervalo entre explicações, mas parte do trabalho de construir e revisar hipóteses.

O primeiro encontro serve para nivelar conceitos. Antes de qualquer modelo, a turma discute aplicações cotidianas de inteligência artificial. Depois separa dado, característica, resultado esperado e regra, e em seguida treinamento, previsão, classificação e agrupamento, até apontar o que seria entregue ao sistema e o que se espera dele. O diagnóstico inicial é individual e acontece no começo desse encontro, antes da discussão, ou na aula anterior, conforme o tempo. Registra ideias de partida, sem antecipar os nomes técnicos que aparecem depois.

Do segundo ao sexto encontro entram os módulos de regressão linear, regressão logística, \textit{k}-NN, árvore de decisão e \textit{k}-means. No segundo, a coleta e a organização da tabela da turma vêm antes da regressão linear. Cada aula retoma elementos das anteriores, para que nenhum modelo seja deixado de lado depois de uma única experiência.

A regressão linear e a regressão logística abrem a parte dos modelos supervisionados. Na linear, a turma reconhece uma tendência, produz uma previsão e a compara com o valor observado, o que traz o gráfico cartesiano, a reta e uma noção inicial de erro. A logística retoma a ideia de prever a partir de exemplos e leva a discussão para a probabilidade e a classificação por limiar. Como as duas aparecem em sequência, a turma compara duas saídas, um valor numérico e uma categoria, e vê que ambas dependem dos dados e das escolhas que estruturam o modelo.

O \textit{k}-NN e a árvore de decisão são duas maneiras de chegar a uma classificação. No \textit{k}-NN, os estudantes observam o papel da distância e do número de vizinhos na votação; na árvore, acompanham regras condicionais e os caminhos até uma decisão. Comparar os dois afasta a ideia de que exista um algoritmo melhor que os outros. Ambos usam dados e regras, cada um à sua maneira, e entregam saídas que precisam ser lidas à luz do problema e dos exemplos disponíveis.

O \textit{k}-means traz uma mudança: os pontos não vêm com categorias dadas de antemão. A atividade volta à pergunta sobre o que se quer produzir com os dados e mostra que formar grupos por proximidade não é escolher uma classe conhecida. Os estudantes levam tempo para atualizar os centroides e observar os pontos se reorganizarem, pois a diferença entre agrupamento e classificação costuma exigir que o professor retome a linguagem dos encontros anteriores. Uma síntese breve registra o que os cinco modelos têm em comum, dados e regras que produzem uma saída, e o que muda conforme a finalidade de cada tarefa.

O sétimo encontro compara os cinco modelos. A turma analisa situações-problema e discute se a tarefa pede prever um valor, escolher uma categoria ou formar grupos, justificando com os dados disponíveis. No oitavo, as duplas organizam uma pequena tabela própria e discutem qual modelo combina com a tarefa. Depois, cada estudante faz um desafio individual de síntese. Precisa reconhecer uma situação, identificar os dados necessários, indicar um tipo de modelo compatível e explicar uma consequência possível de alterar um dado, um parâmetro ou uma regra.

O encontro comparativo não exige domínio de uma taxonomia completa de algoritmos. Recupera diferenças que podem ter ficado dispersas e prepara a transferência para o desafio final. Ao justificar por que uma situação pede previsão, classificação ou agrupamento, o estudante pensa na finalidade da tarefa antes de escolher o nome do modelo. O desafio individual segue essa lógica e verifica se relações trabalhadas em exemplos conhecidos servem a uma situação estruturalmente semelhante, com dados apresentados de outra forma.

O trabalho na plataforma será em duplas, cada uma com um computador do laboratório. As duplas favorecem a conversa sobre hipóteses. As previsões são individuais e depois discutidas em dupla; o desafio final também é individual, para que cada estudante registre a própria compreensão. Se o laboratório estiver indisponível, o professor pode projetar o experimento para a turma e manter os momentos individuais de previsão e explicação.

O detalhamento preliminar dos oito encontros reúne objetivos de aprendizagem e tempo, ações de mediação e questões formativas, e ainda as adaptações de infraestrutura e os limites didáticos de cada módulo. As previsões feitas durante a exploração têm função formativa e não formam corpus individual, porque o site não as armazena. As fontes da análise continuam sendo o diagnóstico inicial, o desafio final e os registros docentes.

A versão do site usada no campo será conferida, antes da intervenção, com os recursos previstos na sequência. Se o cenário da plataforma tiver dados diferentes dos discutidos no roteiro da aula, a diferença será corrigida ou explicada aos estudantes; o mesmo vale para os textos de apoio e as representações matemáticas.

A pesquisa não pode analisar estudantes como se tivessem passado por uma sequência diferente do recurso usado. A conferência é, por isso, condição prática para que sequência, plataforma e instrumentos tenham o mesmo objeto de aprendizagem.

\section{Materiais e limites do produto}

Além da plataforma, a aplicação terá um diagnóstico inicial breve, um desafio final individual e registros docentes sucintos. Não há previsão de uma cartilha extensa para cada aula, porque ela aumentaria a carga de aplicação e de correção sem ser indispensável à investigação. As perguntas do site dão devolutiva imediata durante a aprendizagem. O diagnóstico e o desafio final produzem as evidências que a pesquisa analisa.

O produto não comprova aprendizagem só por existir. Serve para que professor e estudantes discutam dados, parâmetros, cálculos e resultados em sala.

A adequação didática do produto só poderá ser avaliada com a aplicação planejada, as condições reais do campo e as evidências dos estudantes; ele não será apresentado como solução pronta para qualquer escola. Antes da aplicação, uma verificação prática testará a abertura local da plataforma, a leitura dos textos em tela, o funcionamento dos controles e a projeção para a turma toda. Se algum equipamento falhar, a atividade segue com o experimento projetado e pausas para previsão e explicação individual.

A mudança de organização vai para o registro docente, porque a forma de acesso ao recurso pode alterar as interações observadas. Manter uma atividade pedagogicamente equivalente vale mais do que insistir em uma configuração técnica que a escola não consiga sustentar.

A mesma verificação confere se os enunciados se entendem sem que o estudante conheça, antes da aula, o vocabulário técnico. ``Característica'', ``limiar'', ``vizinho'' e ``centroide'' aparecem com explicações situadas e exemplos visíveis. Se uma formulação gerar dúvida repetida, o ajuste do texto ou da mediação entra como aperfeiçoamento do produto, não como falha do estudante. Mantém-se o desafio conceitual; saem apenas os obstáculos de linguagem ou de acesso alheios ao conteúdo.

\section{Alinhamento curricular preliminar}

Uma matriz explicitará o vínculo da sequência com a Computação na Educação Básica, relacionando habilidade curricular, objetivo de aprendizagem, atividade e evidência a observar. A versão preliminar toma como referências do Ensino Médio os fundamentos, as potencialidades e os limites da inteligência artificial, os modelos computacionais simples para simulação e previsão e a análise de dados com princípios de ciência de dados. Elas correspondem às habilidades EM13CO10, EM13CO11 e EM13CO12 do anexo da Resolução CNE/CEB nº 1/2022 \cite{brasil2022computacao}.

A norma não prescreve os cinco modelos como algoritmos. Eles entram como situações didáticas para trabalhar dados, modelos, regras, parâmetros, previsão, classificação, agrupamento e limites de sistemas de IA. A matriz apontará ligações possíveis com a Matemática, sem tirar da Computação o lugar de objeto principal.

A escolha final das habilidades depende da série, do componente curricular, do currículo da rede e das condições confirmadas no campo.
